\documentclass[../../main.tex]{subfiles}
\graphicspath{{../../image/}{../../figure/}}
\begin{document}

\section{补偿紧致性与 div-curl 引理}\label{section:ch17-补偿紧致与divcurl}

补偿紧致性理论致力于对 $L^p$ 空间中赋予了额外结构的有界（弱紧）集进行分类：这种结构达不到（强）紧性，却仍能使某些非线性函数成为弱连续的。下述命题——著名的 div-curl 引理——很好地说明了这一点，它的应用范围惊人地广泛。

\begin{theorem}[div-curl 引理]\label{theorem:ch17-divcurl引理}
设 $\Omega$ 是 $\mathbb{R}^m$ 的开子集，$\{G_j\}$ 与 $\{H_j\}$ 是 $L^2(\Omega;\mathbb{R}^m)$ 中的向量场序列，当 $j\to\infty$ 时分别弱收敛于极限 $\bar G$ 与 $\bar H$。假设 $\{\operatorname{div}G_j\}$ 与 $\{\operatorname{curl}H_j\}$ 均含于 $W^{-1,2}(\Omega)$ 的紧集。则
\begin{align}
G_j\cdot H_j\longrightarrow\bar G\cdot\bar H,\qquad j\to\infty,
\label{align:ch17-1721}
\end{align}
在分布意义下成立。
\end{theorem}

\begin{proof}
只需对 $\Omega$ 有界的情形证明\eqref{align:ch17-1721}即可。此外，由于 $G_j\cdot\bar H\to\bar G\cdot\bar H$，我们可不失一般性地假设 $\bar H=0$。

设 $\Phi_j\in W^{1,2}_0(\Omega;\mathbb{R}^m)\cap W^{2,2}_{\mathrm{loc}}(\Omega;\mathbb{R}^m)$ 表示边值问题 $\Delta\Phi_j=H_j$（在 $\Omega$ 内）、$\Phi_j=0$（在 $\partial\Omega$ 上）的解。则 $\{\Phi_j\}$ 在 $W^{2,2}_{\mathrm{loc}}$ 中弱收敛于零，从而 $\{\operatorname{div}\Phi_j\}$ 在 $W^{1,2}_{\mathrm{loc}}$ 中弱收敛于零。另一方面，由于 $\Delta(\operatorname{curl}\Phi_j)=\operatorname{curl}H_j$，$\{\operatorname{curl}\Phi_j\}$ 在 $W^{1,2}_{\mathrm{loc}}$ 中强收敛于零。

现令
\begin{align}
V_j=H_j-\operatorname{grad}\operatorname{div}\Phi_j
\label{align:ch17-1722}
\end{align}
并注意到，对 $\alpha=1,\ldots,m$，
\begin{align}
V_{j\alpha}=\sum_{\beta=1}^{m}\partial_\beta(\partial_\beta\Phi_{j\alpha}-\partial_\alpha\Phi_{j\beta}),
\label{align:ch17-1723}
\end{align}
于是 $\{V_j\}$ 在 $L^2_{\mathrm{loc}}$ 中强收敛于零。

借助\eqref{align:ch17-1722}，我们得到
\begin{align}
G_j\cdot H_j=G_j\cdot V_j+\operatorname{div}\bigl[(\operatorname{div}\Phi_j)G_j\bigr]-(\operatorname{div}\Phi_j)(\operatorname{div}G_j).
\label{align:ch17-1724}
\end{align}
当 $j\to\infty$ 时，\eqref{align:ch17-1724}右端每一项在分布意义下都趋于零，这就证明了\eqref{align:ch17-1721}。证明完成。
\end{proof}

在应用中，下述技术性结果对于验证\cref{theorem:ch17-divcurl引理}的假设常常很有帮助。

\begin{lemma}\label{lemma:ch17-插值紧致引理}
设 $\Omega$ 是 $\mathbb{R}^m$ 的开子集，$\{\phi_j\}$ 是 $W^{-1,p}(\Omega)$（某 $p>2$）中的有界序列。此外，设 $\phi_j=\chi_j+\psi_j$，其中 $\{\chi_j\}$ 含于 $W^{-1,2}(\Omega)$ 的一个紧集，而 $\{\psi_j\}$ 含于测度空间 $M(\Omega)$ 的一个有界集。则 $\{\phi_j\}$ 含于 $W^{-1,2}(\Omega)$ 的一个紧集。
\end{lemma}

\begin{proof}
考虑 $W^{1,2}_0(\Omega)$ 中（唯一）的函数 $g_j$ 与 $h_j$，它们满足方程
\begin{align}
\Delta g_j=\chi_j,\qquad\Delta h_j=\psi_j.
\label{align:ch17-1725}
\end{align}
由标准椭圆理论，$\{g_j\}$ 含于 $W^{1,2}_0(\Omega)$ 的一个紧集，而 $\{h_j\}$ 含于 $W^{1,q}_0(\Omega)$（$1<q<\frac{m}{m-1}$）的一个紧集。由于 $\phi_j=\Delta(g_j+h_j)$，$\{\phi_j\}$ 含于 $W^{-1,q}(\Omega)$ 的一个紧集。但 $\{\phi_j\}$ 在 $W^{-1,p}(\Omega)$（$p>2$）中有界，于是由 $W^{-1,q}$ 与 $W^{-1,p}$ 之间的插值可知，$\{\phi_j\}$ 含于 $W^{-1,2}(\Omega)$ 的一个紧集。证明完成。
\end{proof}

\end{document}
